\documentclass[11pt]{article}
\usepackage[T1]{fontenc}
\usepackage{lmodern}
\usepackage[margin=1in]{geometry}
\usepackage{amsmath,amssymb,amsthm,mathtools}
\usepackage{booktabs,longtable,array,microtype,xcolor,xurl}
\usepackage[colorlinks=true,linkcolor=blue!40!black,citecolor=blue!40!black,urlcolor=blue!40!black]{hyperref}
\usepackage{bookmark,etoolbox}
\usepackage{fancyhdr}
\pagestyle{fancy}\fancyhf{}
\fancyhead[L]{\small Finite distribution modules and spectral jets}
\fancyhead[R]{\small ProveIt research contribution}
\fancyfoot[C]{\thepage}
\setlength{\headheight}{14pt}
\setlength{\emergencystretch}{2em}
\numberwithin{equation}{section}
\newtheorem{theorem}{Theorem}[section]
\newtheorem{proposition}[theorem]{Proposition}
\newtheorem{lemma}[theorem]{Lemma}
\newtheorem{corollary}[theorem]{Corollary}
\theoremstyle{definition}
\newtheorem{definition}[theorem]{Definition}
\newtheorem{conjecture}[theorem]{Conjecture}
\theoremstyle{remark}
\newtheorem{remark}[theorem]{Remark}
\newtheorem{example}[theorem]{Example}
\DeclareMathOperator{\Li}{Li}
\DeclareMathOperator{\rank}{rank}
\DeclareMathOperator{\ord}{ord}
\DeclareMathOperator{\rad}{rad}
\DeclareMathOperator{\Ree}{Re}
\DeclareMathOperator{\Imm}{Im}
\DeclareMathOperator{\cond}{cond}
\newcommand{\Q}{\mathbb Q}
\newcommand{\C}{\mathbb C}
\newcommand{\Z}{\mathbb Z}
\newcommand{\N}{\mathbb N}
\newcommand{\ee}{\mathrm e}
\newcommand{\ii}{\mathrm i}
\newcommand{\U}{\mathcal U}
\newcommand{\V}{\mathcal V}
\newcommand{\D}{\mathcal D}
\newcommand{\cP}{\mathcal P}
\newcommand{\wt}{\widetilde}
\newcommand{\code}[1]{\texttt{\small #1}}
\newcommand{\Liord}[2]{\left.\partial_s^{#1}\Li_s(#2)\right|_{s=1}}
\title{\textbf{Finite Distribution Modules, Spectral Jets,\\ and Cyclotomic Trace Identities}\\[8pt]
\large Uniform rank proofs, explicit normal forms, exceptional zeros,\\
and an exact audit of a Gaussian weight-five conjecture}
\author{Research contribution prepared for the ProveIt project\\
\small OpenAI-assisted mathematical research; supplied for independent review}
\date{9 October 2026}
\hypersetup{pdftitle={Finite Distribution Modules, Spectral Jets, and Cyclotomic Trace Identities},pdfauthor={Research contribution prepared for ProveIt}}
\begin{document}
\maketitle
\begin{abstract}
The ProveIt polylogarithm manuscript records finite experiments suggesting
that rational distribution relations leave $\varphi(q)-1$ coordinates at
level $q$, and $\varphi(q)/2-1$ after suitable reflection relations. We give
a uniform proof for an explicitly defined presentation, simultaneously for
all independent prime weights and every finite spectral-jet order. Over a
character splitting field, the universal weighted quotient is a free
polynomial module of rank $\varphi(q)$. Its conductor-indexed basis remains
valid at resonant weights, whereas the basis consisting of primitive-grid
points can fail. We factor the determinant of that change of coordinates,
compute its local jet defects, and give rational primitive-grid normal
forms with exact support, sign, and denominator information.

The resulting calculus yields all-orders Hurwitz--Stieltjes reduction
identities, including the pole correction for underived Stieltjes
constants. Applying the same module to roots-of-unity polylogarithms gives
character trace identities and exact orders of vanishing at spectral
order one; for example, the primitive level-$30$ trace has second
order derivative $-2\log2\log3\log5$. Separately, a rational certificate
shows that the manuscript's $S_4$ conjecture is not a consequence of a
specified $96$-row weight-five double-shuffle presentation, without
asserting that the conjecture is false. An audit also corrects the status
of the cubic log-gamma moment using the published work of
Bailey--Borwein--Borwein. Complete proofs, exact regression programs,
replayable certificates, numerical diagnostics, and integration notes
accompany the article.
\end{abstract}
\noindent\textbf{Scope of claims.} All rank statements concern explicitly
presented formal modules. They do not assert linear independence of
transcendental values. No proof-assistant verification or global priority
claim is made. Classical antecedents and the remaining conjecture are
identified explicitly.
\newpage
\begingroup
\small
\makeatletter
\patchcmd{\l@section}{1.0em}{0.4em}{}{}
\makeatother
\tableofcontents
\endgroup
\newpage

\section{Research target, provenance, and mathematical context}\label{sec:scope}

\subsection{The question resolved here}
The starting point is Chapters 4, 7, 8, and 10 of the ProveIt manuscript
\emph{Polylogarithms and their Arithmetic Bridges}~\cite{repo}.
Chapter 8, Observation \code{tower:thm:rank}, reports exact finite
experiments with $3\le q\le30$ and three parameter-derivative orders:
\begin{equation}\label{eq:source-rank}
(q-1)-\rank(\text{distribution rows})=\varphi(q)-1.
\end{equation}
Chapter 7 gives the analogous count $\varphi(q)/2-1$ after distribution
and first-Stieltjes reflection. The manuscript correctly distinguishes
these formal calculations from numerical independence and explicitly
asks for a proof uniform in the denominator and derivative order.

The historical matrices are not included in the cited manuscript
snapshot. Consequently we do not claim to reconstruct their exact row
ordering or to certify unseen computations. Instead, we define the
canonical complete divisor-row presentation from the stated distribution
law, prove its rank for every $q$, and give code generating every row.
This resolves the uniform mathematical problem underlying the observation
and supplies a precise replacement for its experimental formulation.

The stronger result is useful: the rank does not depend on the
multiplicative weights at all. Apparent failures at exceptional spectral
parameters are failures of a chosen coordinate basis, not jumps in the
rank of the complete relation module. A determinant factorization makes
this distinction quantitative.

\subsection{Classical antecedents and the contribution's status}
Universal distributions are classical. In particular, the ordinary
specialization of our module, with all prime weights equal to one,
is a finite-level universal ordinary distribution. Ouyang's
Proposition 3.1 records its integral freeness and rank
$\varphi(q)$, citing the work of Kubert and Anderson~\cite{ouyang,kubert}.
We do not present that specialization as a new theorem.

The contribution developed here is an explicit polynomial-weight
presentation with conductor normal forms, its determinant and jet
specializations, and the resulting integration-ready proof of the
manuscript's rank problem. The polylogarithm trace formulas derive from
classical distribution and Gauss-sum identities; their concrete
order-derivative evaluations are supplied as consequences, without a
claim that each has not appeared elsewhere. The exact $S_4$ row-space
obstruction concerns a newly specified finite presentation, rather than
all possible multiple-polylogarithm identities.

Standard analytic conventions for Hurwitz zeta, polylogarithms, and
Dirichlet $L$-functions are as in the corresponding DLMF sections
\cite{dlmf-hurwitz,dlmf-polylog,dlmf-dirichlet}. Every additional algebraic
rank assertion used here is proved below rather than inferred from a
numerical relation search.

\subsection{Notation and endpoint convention}
For $q\ge2$, let
\[
 X_q=(q^{-1}\Z)/\Z,\qquad U_m=(\Z/m\Z)^\times,
 \qquad P(q)=\{p:p\text{ prime},\ p\mid q\}.
\]
We use $U_1$ as a one-element group and put $\varphi(1)=1$.
The point $0\in X_q$ is a \emph{formal} point. Under Hurwitz evaluation it
means the endpoint $1$, never the undefined expression $\zeta(s,0)$:
\[
 e_{a/m}\longmapsto\zeta(s,a/m)\quad(1\le a<m),
 \qquad e_0\longmapsto\zeta(s,1).
\]
Under cyclotomic evaluation it means $\Li_s(1)=\zeta(s)$, initially
away from the pole. This convention is essential when a distribution
fiber passes through zero.

Every character of $U_q$ has a unique primitive inducing character of
conductor $f\mid q$. A symbol $\chi$ will denote that primitive character
and, where specified, its inflation to $U_m$ for $f\mid m\mid q$.
The principal character has conductor one. Character values at an
integer prime to $f$ are roots of unity; primes dividing $f$ are treated
separately, so no ambiguous evaluation at a nonunit is needed.

\section{The universal weighted distribution module}\label{sec:universal}

Fix a characteristic-zero field $K$ containing all values of characters
of $U_q$, and put
\[
 R=K[A_p:p\in P(q)],\qquad
 \V_q=\bigoplus_{x\in X_q}R e_x.
\]
The $A_p$ are algebraically independent indeterminates. For each
$p\mid q$ and $x\in (p/q)\Z/\Z$, define a prime distribution row
\begin{equation}\label{eq:prime-row}
 r_{p,x}=\sum_{\substack{y\in X_q\\py=x}}e_y-A_p e_x.
\end{equation}
Let $\mathcal R_q$ be their $R$-span and let
$\D_q=\V_q/\mathcal R_q$.

For $d\mid q$ write $A_d=\prod_{p\mid d}A_p^{v_p(d)}$.
The corresponding divisor row is
\begin{equation}\label{eq:divisor-row}
 r_{d,x}=\sum_{dy=x}e_y-A_d e_x,
 \qquad x\in(d/q)\Z/\Z.
\end{equation}
The multiplicativity of $A_d$ is part of the definition; arbitrary
unrelated weights on composite divisors would be a different problem.

\begin{lemma}[Prime rows suffice]\label{lem:prime-suffice}
Every row \eqref{eq:divisor-row} lies in $\mathcal R_q$.
\end{lemma}
\begin{proof}
If $d=ab\mid q$, then for $x\in(d/q)\Z/\Z$ every solution of $ay=x$
lies in $(b/q)\Z/\Z$, and
\[
 r_{ab,x}=\sum_{ay=x}r_{b,y}+A_b r_{a,x}.
\]
Each $z$ with $abz=x$ occurs exactly once in the double sum, and the
intermediate $e_y$ terms cancel. Induction on the number of prime factors
of $d$, counted with multiplicity, proves the assertion. No division by
any weight is used.
\end{proof}

For $f=\cond(\chi)$ and $f\mid m\mid q$, set
\begin{equation}\label{eq:v-character}
 v_{m,\chi}=\sum_{a\in U_m}\chi(a)e_{a/m},
 \qquad v_{1,\mathbf1}=e_0.
\end{equation}
Define the polynomial
\begin{equation}\label{eq:E-polynomial}
 E_{m,\chi}=
 \prod_{p\mid f}A_p^{v_p(m)-v_p(f)}
 \prod_{\substack{p\mid m\\p\nmid f}}
 A_p^{v_p(m)-1}\bigl(A_p-\chi(p)\bigr).
\end{equation}
In particular $E_{f,\chi}=1$.

\begin{theorem}[Universal conductor normal form]\label{thm:universal}
The quotient $\D_q$ is a free $R$-module of rank $\varphi(q)$.
One basis consists of
\[
 b_\chi=[v_{f,\chi}],\qquad f=\cond(\chi),\quad \chi\bmod q.
\]
In this basis,
\begin{equation}\label{eq:conductor-normal}
 [v_{m,\chi}]=E_{m,\chi}b_\chi\qquad(f\mid m\mid q).
\end{equation}
The relation module is free of rank $q-\varphi(q)$ and is a direct
summand of $\V_q$. All these statements survive arbitrary specialization
of the weights into any $K$-algebra, including $A_p=0$ and $A_p=1$.
\end{theorem}

\subsection{Proof: identify the complete relation module}
It is not enough to exhibit character identities satisfied by evaluated
special functions: that would give only one direction of a rank argument.
We therefore identify the image of the entire relation map.

The points of $X_q$ split by exact denominator. Hence
\begin{equation}\label{eq:denominator-split}
 \V_q=\bigoplus_{m\mid q}R[U_m].
\end{equation}
The reduction map $U_q\to U_m$ is surjective. Character orthogonality
therefore gives an invertible Fourier change of basis on each summand.
A character of conductor $f$ occurs exactly once in the summand indexed
by $m$ when $f\mid m$, and does not occur otherwise. Thus the $\chi$
block of \eqref{eq:denominator-split} has basis
\[
 v_{m,\chi},\qquad f\mid m\mid q.
\]

Fix $p\mid q$. The source labels $x$ of the rows $r_{p,x}$ themselves
form a grid of level $q/p$. Fourier-transforming separately on its
exact-denominator strata gives all transformed rows, not a selected
subset of them. For $f\mid m$ and $pm\mid q$, the transformed row is
\begin{equation}\label{eq:edge-relations}
 v_{pm,\chi}-\lambda_{p,m,\chi}v_{m,\chi},
 \qquad
 \lambda_{p,m,\chi}=
 \begin{cases}
 A_p,&p\mid m,\\
 A_p-\chi(p),&p\nmid m.
 \end{cases}
\end{equation}
To check this directly, weight the row at $x=a/m$ by $\chi(a)$ and sum
over $U_m$. If $p\mid m$, all $p$ preimages have exact denominator $pm$.
If $p\nmid m$, one preimage has denominator $m$: it is $b/m$ with
$pb\equiv a\pmod m$, and its coefficient is
$\chi(a)=\chi(p)\chi(b)$. All other preimages have denominator $pm$.
This proves \eqref{eq:edge-relations}. Since the Fourier transforms of
the source and target are invertible, these edges span precisely the
original relation module in character coordinates.

Starting at $m=f$, multiply the edge factors while adjoining prime
factors of $m/f$. A prime already dividing $f$ always contributes $A_p$.
A prime not dividing $f$ contributes $A_p-\chi(p)$ on its first
appearance and $A_p$ on every later appearance. The resulting product
is exactly \eqref{eq:E-polynomial}, independently of the order of
adjoining primes. Therefore all $v_{m,\chi}$ reduce to
$E_{m,\chi}v_{f,\chi}$.

There are no additional relations on $v_{f,\chi}$. Indeed, the map from
the free $\chi$ block to $R$ defined by
\[
 v_{m,\chi}\longmapsto E_{m,\chi}
\]
annihilates every edge and sends $v_{f,\chi}$ to one. It induces an
inverse to $R\to\D_{q,\chi}$, $1\mapsto[v_{f,\chi}]$. Moreover, its
kernel has the explicit monic basis
\begin{equation}\label{eq:monic-basis}
 v_{m,\chi}-E_{m,\chi}v_{f,\chi},\qquad
 f\mid m\mid q,\ m\ne f.
\end{equation}
These vectors are independent by their distinct non-base coordinates.
There is one quotient coordinate for each character of $U_q$, hence
$\varphi(q)$ in all, and the complementary kernel has rank
$q-\varphi(q)$. This is a split decomposition over the polynomial ring
itself, so any base change preserves it.\qed

\begin{corollary}[Pointwise polynomial normal form]\label{cor:point-normal}
For a reduced point $a/m$ with $m\mid q$,
\begin{equation}\label{eq:point-normal}
 [e_{a/m}]=\frac1{\varphi(m)}
 \sum_{\chi\bmod m}\overline{\chi(a)}\,E_{m,\chi}b_\chi.
\end{equation}
The sum runs over characters of $U_m$, identified with their primitive
inducing characters and then inflated to $U_q$.
\end{corollary}
\begin{proof}
Apply inverse character Fourier transformation to
\eqref{eq:v-character}, then use \eqref{eq:conductor-normal}.
\end{proof}

\subsection{A replayable symbolic reduction certificate}
Choose a chain
$f=m_0\mid m_1\mid\cdots\mid m_t=m$ in which each ratio is prime.
Let $\lambda_j$ be its edge factor. The exact identity
\begin{equation}\label{eq:telescoping-certificate}
 v_{m,\chi}-E_{m,\chi}v_{f,\chi}
 =\sum_{j=1}^t\left(\prod_{h=j+1}^t\lambda_h\right)
 \bigl(v_{m_j,\chi}-\lambda_jv_{m_{j-1},\chi}\bigr)
\end{equation}
is a finite row-combination certificate. Its verification is polynomial
cancellation. In particular, a certificate does not become invalid when
one or several factors vanish. This is the practical advantage of
conductor coordinates over a generic matrix inverse.

\section{Anchoring, reflection, and uniform jet ranks}\label{sec:jets}

\subsection{The precise replacement for the finite rank observation}
For actual evaluations, the endpoint value and lower-order terms can be
regarded as known. In a homogeneous difference system this amounts to
setting the corresponding formal coordinates equal to zero.

\begin{corollary}[Anchored rank]\label{cor:anchored}
At every specialization of the weights, the quotient by the distribution
relations and $e_0=0$ has dimension $\varphi(q)-1$. Equivalently, after
deleting the endpoint column, the prime or complete divisor-row matrix
on the remaining $q-1$ columns has rank $q-\varphi(q)$.
\end{corollary}
\begin{proof}
The endpoint is $b_{\mathbf1}$, the basis vector of conductor one.
It is nonzero and independent of every other conductor coordinate.
Adding its row increases the full relation rank by one. Eliminating
that endpoint column leaves rank $q-\varphi(q)$ on $q-1$ columns.
\end{proof}

This statement applies, in particular, to the homogeneous highest-order
part of the manuscript's $\gamma_1^{(k)}$ distribution equations, where
$A_p=p^{k+1}$. It holds for every $q\ge2$ and $k\ge1$, not only the
finite experimental range. Known inhomogeneous right sides do not alter
this rank provided the system is consistent; the analytic evaluations
supply consistency.

The reflection operator $J(e_x)=e_{-x}$ preserves the relation module.
Since $Jv_{m,\chi}=\chi(-1)v_{m,\chi}$, its action is diagonal on the
conductor basis.

\begin{corollary}[Parity counts]\label{cor:parity}
For $q>2$, the unanchored even and odd quotients have dimensions
$\varphi(q)/2$ each. After anchoring the endpoint, their dimensions are
respectively $\varphi(q)/2-1$ and $\varphi(q)/2$.
Thus imposing known odd reflection differences leaves exactly
$\varphi(q)/2-1$ formal homogeneous coordinates.
\end{corollary}
\begin{proof}
The element $-1\in U_q$ has order two and is not the identity. Exactly
half the characters take value $1$ there and half take value $-1$.
The principal character is even, so anchoring removes one even
coordinate. For $q=2$ the separate counts are zero even coordinates
after anchoring and zero odd coordinates.
\end{proof}

This proves the canonical formal version of the Chapter 7 count. It does
not say that $\varphi(q)/2-1$ numerical Stieltjes values are independent:
additional identities may impose further relations after evaluation.

\subsection{All finite jets by base change}
Let $F$ be a field extending $K$, and consider the truncated ring
\[
 T_N=F[t]/(t^{N+1}),\qquad N\ge0.
\]
Choose arbitrary truncated prime weights $A_p(t)\in T_N$. Of particular
interest are
\begin{equation}\label{eq:spectral-weight}
 A_p(t)=p^{s_0}\exp(t\log p)\pmod{t^{N+1}}
\end{equation}
for Hurwitz zeta, and $p^{1-s_0}\exp(-t\log p)$ for polylogarithms.
The coefficient field may be $\C$, or a field containing the displayed
constants; no algebraic independence of actual logarithms is assumed.

\begin{theorem}[Uniform finite-jet ranks]\label{thm:jet-rank}
The truncated distribution quotient is a free $T_N$-module of rank
$\varphi(q)$ for every choice of $A_p(t)$. Its dimension over $F$ is
$(N+1)\varphi(q)$, and the full coefficient-expanded relation matrix
has rank
\begin{equation}\label{eq:jet-rank}
 (N+1)\bigl(q-\varphi(q)\bigr).
\end{equation}
After fixing all endpoint coefficients its dimension is
$(N+1)(\varphi(q)-1)$. The anchored even and odd dimensions for $q>2$
are $(N+1)(\varphi(q)/2-1)$ and $(N+1)\varphi(q)/2$.
\end{theorem}
\begin{proof}
Apply the homomorphism $R\to T_N$ that sends each indeterminate to the
chosen truncated weight. The split decomposition proved in
Theorem~\ref{thm:universal} remains split under this base change.
Expanding in the basis $1,t,\ldots,t^N$ multiplies every free rank by
$N+1$. Anchoring and reflection act on the same conductor coordinates
as before. This argument is stronger than an appeal to a block-triangular
matrix with a possibly singular diagonal: it applies at every weight.
\end{proof}

If all lower jets are already fixed and only the highest jet is unknown,
the homogeneous matrix is the order-zero distribution matrix. Its
anchored residual dimension is therefore $\varphi(q)-1$. The full
stacked system and this one-layer system are different matrices, and
\eqref{eq:jet-rank} should not be confused with the one-layer rank.

\section{The primitive-grid determinant and its exceptional locus}\label{sec:det}

Conductor coordinates are polynomial and globally valid, but it is often
more convenient to use the $\varphi(q)$ primitive-grid points
$e_{u/q}$, $u\in U_q$. Fourier-transforming those points gives
$v_{q,\chi}$. By Theorem~\ref{thm:universal}, the map from their formal
span into $\D_q$ is diagonal in character coordinates:
\begin{equation}\label{eq:primitive-diagonal}
 v_{q,\chi}\longmapsto E_{q,\chi}b_\chi.
\end{equation}
Its determinant in these specified Fourier and conductor bases is
$\Delta_q=\prod_{\chi\bmod q}E_{q,\chi}$.

For $p^e\Vert q$, write
\[
 M_p=q/p^e,\qquad h_p=\ord_{M_p}(p),
\]
with $h_p=1$ when $M_p=1$.

\begin{theorem}[Determinant factorization]\label{thm:determinant}
One has the polynomial identity
\begin{equation}\label{eq:determinant}
 \boxed{\displaystyle
 \Delta_q(\mathbf A)=
 \prod_{p^e\Vert q}
 A_p^{\varphi(M_p)(p^{e-1}-1)}
 \bigl(A_p^{h_p}-1\bigr)^{\varphi(M_p)/h_p}.}
\end{equation}
Thus primitive-grid coordinates form a basis precisely when all the
factors in \eqref{eq:determinant} are invertible in a field
specialization. Their failure does not change the rank of $\D_q$.
\end{theorem}
\begin{proof}
Fix $p^e\Vert q$ and inspect the $p$ factor of
$\prod_\chi E_{q,\chi}$. Characters whose conductor is not divisible by
$p$ are exactly the characters of $U_{M_p}$. Their additional factors
are $A_p-\chi(p)$. In the character group of $U_{M_p}$, the values
$\chi(p)$ run through every $h_p$th root of unity with multiplicity
$\varphi(M_p)/h_p$. This follows by duality from the fact that $p$ has
order $h_p$ in $U_{M_p}$. Their product is therefore
$(A_p^{h_p}-1)^{\varphi(M_p)/h_p}$.

It remains to count the pure powers of $A_p$. If $a=v_p(f)$, the
exponent contributed by one character is $e-\max(1,a)$. Writing this
as a sum of indicators gives
\[
 \sum_{\chi\bmod q}\bigl(e-\max(1,v_p(f_\chi))\bigr)
 =\sum_{j=1}^{e-1}\#\{\chi:v_p(f_\chi)\le j\}.
\]
The count in the last sum is $\varphi(p^j)\varphi(M_p)$, since these
are exactly the characters induced from level $p^jM_p$. Hence the
exponent is
$\varphi(M_p)\sum_{j=1}^{e-1}\varphi(p^j)
=\varphi(M_p)(p^{e-1}-1)$. Multiply over the primes.
\end{proof}

\begin{example}[Level twelve]
Let $\chi_3$ and $\chi_4$ be the nonprincipal quadratic characters of
conductors three and four, and let $\chi_{12}=\chi_3\chi_4$.
Their conductor generators are
\begin{align*}
 b_3&=e_{1/3}-e_{2/3},&
 b_4&=e_{1/4}-e_{3/4},\\
 b_{12}&=e_{1/12}-e_{5/12}-e_{7/12}+e_{11/12}.
\end{align*}
For $r\in U_{12}$, the normal form reads
\begin{align}\label{eq:q12-normal}
 [e_{r/12}]=\frac14\bigl[&A_2(A_2-1)(A_3-1)e_0
 +\chi_3(r)A_2(A_2+1)b_3\notag\\
 &+\chi_4(r)(A_3+1)b_4+\chi_{12}(r)b_{12}\bigr].
\end{align}
Consequently
\[
 \Delta_{12}=A_2^2(A_2^2-1)(A_3^2-1).
\]
At $A_2=A_3=1$, the primitive-grid span misses the principal conductor
coordinate, even though the full quotient still has rank four.
\end{example}

\subsection{Local jet defects of the primitive-grid map}
Specialize to $A_p(s)=p^s$, and let $s=s_0+t$. All exponential factors
are nonzero. A factor $p^s-\chi(p)$ vanishes precisely when
$p^{s_0}=\chi(p)$, and any such zero is simple because its derivative
is $p^{s_0}\log p\ne0$.

\begin{theorem}[Local Smith factors and primitive-jet ranks]\label{thm:smith}
Define
\begin{equation}\label{eq:rho-general}
 \rho_\chi(s_0)=\#\{p\mid q:p\nmid f_\chi,\ p^{s_0}=\chi(p)\}.
\end{equation}
Over $\C[[t]]$, the diagonal factors of the primitive-grid map are
units times $t^{\rho_\chi(s_0)}$. Consequently its rank on order-$N$
truncated jets, as a complex linear map, is
\begin{equation}\label{eq:primitive-jet-rank}
 \sum_{\chi\bmod q}\max\{N+1-\rho_\chi(s_0),0\}.
\end{equation}
At $s_0=0$, this is determined by
$\rho_\chi(0)=\#\{p\mid q:p\nmid f_\chi,\ \chi(p)=1\}$.
\end{theorem}
\begin{proof}
Factor $E_{q,\chi}(s_0+t)$ into its simple vanishing factors and a unit.
The primitive Fourier matrix is constant and invertible, so these are
also the local invariant factors, up to permutation and units.
Multiplication by $t^r$ on $\C[t]/(t^{N+1})$ has rank
$\max(N+1-r,0)$. Summing over the diagonal coordinates proves the
formula. This is the rank of a \emph{coordinate inclusion map}, not
the rank of the distribution relation matrix.
\end{proof}

\begin{corollary}[Multiple-prime resonances are central]\label{cor:central}
If $s_0\ne0$, at most one prime can satisfy
$p^{s_0h_p}=1$. At a nonzero resonant point for that prime $p$, the
primitive-grid rank defect is $\varphi(M_p)/h_p$, and every affected
local factor has order one. The order-$N$ inclusion rank is then
$(N+1)\varphi(q)-\varphi(M_p)/h_p$ for all $N\ge0$.
\end{corollary}
\begin{proof}
A resonance forces $s_0$ to be purely imaginary. Resonances at distinct
primes would make $\log p/\log\ell$ rational, because each phase is a
rational multiple of $2\pi$. Unique prime factorization rules this out
unless $s_0=0$. For the one resonant prime, the multiplicity count in
Theorem~\ref{thm:determinant} gives the number of affected characters.
Apply Theorem~\ref{thm:smith}.
\end{proof}

\begin{example}[Two exact jet-rank profiles]
At level $12$ and $s_0=0$, the principal character has $\rho=2$ and the
other three characters have $\rho=0$. For $N=0,1,2$, the primitive-jet
ranks are $3,6,10$, while the full quotient dimensions are $4,8,12$.
At level $21$, the principal character has $\rho=2$, the character
induced from conductor three has $\rho=1$, and the other ten
characters have $\rho=0$. The inclusion ranks are $10,21,33$, compared
with full quotient dimensions $12,24,36$. Indeed,
\[
 \Delta_{21}=(A_3^6-1)(A_7-1)^2.
\]
Exact character-phase certificates for these examples are supplied.
\end{example}

\section{Explicit primitive-grid reduction coefficients}\label{sec:primitive}

Away from the exceptional locus, the conductor normal form can be
inverted to express all grid points through the primitive ones. For
Hurwitz weights there is a particularly concrete answer that avoids
character-field arithmetic.

Let $\mathcal S_q$ be the multiplicative semigroup of positive integers
whose prime factors divide $q$, including $1$. For $\Ree s>0$, define
\begin{equation}\label{eq:C-smooth}
 C_{r,u}(s)=\sum_{\substack{d\in\mathcal S_q\\du\equiv r\; (q)}}d^{-s},
 \qquad 0\le r<q,\quad u\in U_q.
\end{equation}
The series is absolutely convergent, since
$\sum_{d\in\mathcal S_q}d^{-\sigma}
=\prod_{p\mid q}(1-p^{-\sigma})^{-1}$ for $\sigma>0$.

\begin{theorem}[Primitive-grid extension]\label{thm:primitive}
For $\Ree s>0$, the $q\times\varphi(q)$ matrix $C(s)$ satisfies every
prime and divisor distribution relation with weights $p^s$, and its
primitive rows form the identity matrix. It is the unique extension
matrix from primitive-grid data. In particular,
\begin{equation}\label{eq:Hurwitz-C}
 \zeta(s,r/q)=\sum_{u\in U_q}C_{r,u}(s)\zeta(s,u/q)
\end{equation}
with the endpoint convention for $r=0$, initially for $\Ree s>1$ and
then meromorphically wherever the coefficients are defined.
\end{theorem}
\begin{proof}
For a primitive row $r$, the congruence $du\equiv r\pmod q$ forces
$d$ to be prime to $q$, hence $d=1$. This gives the identity submatrix.
For a prime distribution row indexed by $a\equiv0\pmod p$,
\[
 \sum_{pr\equiv a\;(q)}C_{r,u}(s)
 =\sum_{\substack{d\in\mathcal S_q\\pdu\equiv a\;(q)}}d^{-s}
 =p^s\sum_{\substack{d'\in\mathcal S_q\\d'u\equiv a\;(q)}}(d')^{-s}.
\]
In the last step $d'=pd$; conversely every $d'$ in the last sum is
divisible by $p$, because $u$ is a unit and $a$ is divisible by $p$.
The last expression is $p^sC_{a,u}(s)$. Prime sufficiency gives all
divisor rows. Theorem~\ref{thm:universal} and the identity submatrix
show that these $\varphi(q)$ independent columns exhaust the solution
space. Hurwitz zeta satisfies the same distribution equations by
regrouping its absolutely convergent defining series, proving
\eqref{eq:Hurwitz-C}.
\end{proof}

\subsection{A finite formula, including negative spectral parameters}
For $0\le r<q$, let
\[
 g=\gcd(r,q),\quad m=q/g,\quad a=r/g,
 \quad P_r=\{p\mid q:p\nmid m\}.
\]
When $r=0$, this means $g=q$, $m=1$, $a=0$.
For $p\in P_r$, put $h_{p,m}=\ord_m(p)$, with order one modulo one.

\begin{theorem}[Finite residue formula]\label{thm:finite-C}
The meromorphic continuation of \eqref{eq:C-smooth} is
\begin{equation}\label{eq:C-finite}
 C_{r,u}(s)=g^{-s}
 \frac{\displaystyle
 \sum_{\substack{0\le b_p<h_{p,m}\;(p\in P_r)\\
 u\prod_{p\in P_r}p^{b_p}\equiv a\;(m)}}
 \prod_{p\in P_r}p^{-s b_p}}
 {\displaystyle\prod_{p\in P_r}(1-p^{-s h_{p,m}})}.
\end{equation}
Empty products have their usual value one. In particular,
\begin{equation}\label{eq:C-endpoint}
 C_{0,u}(s)=q^{-s}\prod_{p\mid q}(1-p^{-s})^{-1},
\end{equation}
independently of $u$.
\end{theorem}
\begin{proof}
The congruence in \eqref{eq:C-smooth} forces the valuation of $d$ at
an unsaturated prime, one dividing $m$, to be exactly $v_p(g)$.
At a saturated prime, one not dividing $m$, it is at least $v_p(g)$.
Thus $d=g\prod_{p\in P_r}p^{c_p}$ with $c_p\ge0$, and the remaining
congruence is $u\prod p^{c_p}\equiv a\pmod m$.
Write $c_p=b_p+h_{p,m}j_p$ and sum the independent geometric series in
$j_p$. This proves the formula in $\Ree s>0$ and hence its meromorphic
continuation. The endpoint has trivial congruence modulo one.
\end{proof}

Formula \eqref{eq:C-finite}, \emph{not} the divergent smooth-number
series, is the definition used at negative real $s$. It has no pole at
any nonzero real $s$, and the continued distribution identities remain
valid there. At a zero of a displayed denominator, one should use the
polynomial conductor normal form rather than cancel poles numerically.
Some individual coefficients may have removable singularities, but
Theorem~\ref{thm:determinant} gives the criterion for the entire
primitive-grid coordinate system.

\begin{corollary}[Exact support and signs]\label{cor:support}
Let $H_r$ be the subgroup of $U_m$ generated by the classes of the
primes in $P_r$. For every nonzero real $s$,
\[
 C_{r,u}(s)\ne0\quad\Longleftrightarrow\quad au^{-1}\in H_r.
\]
The number of supported primitive coordinates is
\begin{equation}\label{eq:support-count}
 \frac{\varphi(q)}{\varphi(m)}|H_r|.
\end{equation}
All supported coefficients are positive for $s>0$. For $s<0$, they
all have sign $(-1)^{|P_r|}$.
\end{corollary}
\begin{proof}
The numerator of \eqref{eq:C-finite} is a sum of positive real terms
and is nonempty exactly under the subgroup condition. The reduction
map $U_q\to U_m$ has fibers of cardinality
$\varphi(q)/\varphi(m)$, giving \eqref{eq:support-count}.
For $s>0$ every denominator factor is positive. For $s<0$ each of the
$|P_r|$ denominator factors is negative, while the numerator and
$g^{-s}$ are positive.
\end{proof}

\begin{corollary}[Complete monotonicity and normalization]\label{cor:monotone}
For $s>0$ and $j\ge0$, put
\begin{equation}\label{eq:W-definition}
 W_{r,u}^{[j]}(s)=(-1)^j\frac{d^j}{ds^j}C_{r,u}(s)
 =\sum_{\substack{d\in\mathcal S_q\\du\equiv r\;(q)}}
 \frac{(\log d)^j}{d^s}.
\end{equation}
These quantities are nonnegative. Moreover,
$\sum_u C_{r,u}(1)=1$ for every $r$, and
$0\le C_{r,u}(k)\le1$ for integer $k\ge1$.
\end{corollary}
\begin{proof}
The logarithmically weighted semigroup series converges locally
uniformly in $\Ree s>0$, so termwise differentiation is valid.
The constant grid vector satisfies the distribution equations at
$s=1$, and its primitive entries are all one. Uniqueness of extension
proves the row-sum identity. Monotonicity gives
$C_{r,u}(k)\le C_{r,u}(1)\le1$.
\end{proof}

\begin{example}[Two reductions at denominator six]\label{ex:q6}
The nontrivial coefficients for $r=2$ are
\[
 C_{2,1}(s)=\frac{2^s}{2^{2s}-1},\qquad
 C_{2,5}(s)=\frac1{2^{2s}-1}.
\]
Thus, meromorphically,
\begin{align}
 \zeta(s,1/3)&=
 \frac{2^s\zeta(s,1/6)+\zeta(s,5/6)}{2^{2s}-1},\label{eq:q6-one}\\
 \zeta(s,1/2)&=
 \frac{\zeta(s,1/6)+\zeta(s,5/6)}{3^s-1}.
 \label{eq:q6-two}
\end{align}
At $s=1$, the coefficients in \eqref{eq:q6-one} are $2/3$ and $1/3$.
At $s=-1$ they are $-2/3$ and $-4/3$, illustrating the sign theorem
without using a divergent series. The poles at $s=0$ reflect a
singular primitive basis, not a failure of the underlying distribution
relations.
\end{example}

\subsection{Denominator certificates and construction cost}
Let $M_p$ also denote, in this paragraph only, the pushforward matrix
on grid vectors defined by
\[
 (M_pf)_r=\sum_{pu\equiv r\;(q)}f_u.
\]
To avoid confusion with the integer complement in
Section~\ref{sec:det}, write that complement as $q_p=q/p^e$ here and
put $h_p=\ord_{q_p}(p)$. Multiplication of residues gives
$M_p^{e+h_p}=M_p^e$. For $z=p^{-s}$,
\begin{equation}\label{eq:operator-inverse}
 (I-zM_p)^{-1}
 =\sum_{j=0}^{e-1}z^jM_p^j
 +\frac{z^e}{1-z^{h_p}}\sum_{j=0}^{h_p-1}z^jM_p^{e+j}.
\end{equation}
This follows first from a geometric series for $|z|<1$, then as a
rational matrix identity. The commuting product of these inverses,
applied to vectors supported on $U_q$, is the extension $C(s)$.

\begin{proposition}[A uniform rational denominator]\label{prop:denominator}
For integer $k\ge1$, every reduced denominator of $C_{r,u}(k)$ divides
\begin{equation}\label{eq:denominator-bound}
 D_q(k)=\prod_{p^e\Vert q}
 p^{k(e-1)}\bigl(p^{kh_p}-1\bigr).
\end{equation}
Its bit length is $O(kq\log q)$. The coefficients of the raw derivative
$C_{r,u}^{(j)}(k)$, viewed as a polynomial in the symbols $\log p$,
have denominators dividing $D_q(k)^{j+1}$.
\end{proposition}
\begin{proof}
Every scalar coefficient on the right of
\eqref{eq:operator-inverse} has denominator dividing
$p^{k(e-1)}(p^{kh_p}-1)$. Multiplication over primes proves the first
claim. Since $h_p\le q$ and $\sum_{p\mid q}\log p\le\log q$, its
logarithm is at most a constant times $kq\log q$.
The matrix product can be written as a ratio of integer polynomials in
$p^s$, with denominator obtained from \eqref{eq:denominator-bound} by
replacing $k$ with $s$. Repeated differentiation introduces polynomial
factors in $\log p$ and raises the denominator power by at most one
each time. For normalized Taylor coefficients, an additional factor
$j!$ may be needed.
\end{proof}

The implementation computes \eqref{eq:C-finite} by convolution on
residue classes modulo $m$, separately for each $m\mid q$. At a fixed
$m$, each prime update takes at most $\varphi(m)^2$ rational arithmetic
operations. A complete matrix can therefore be constructed in
\[
 O\!\left(\omega(q)\sum_{m\mid q}\varphi(m)^2+q\varphi(q)\right)
 =O\bigl(\omega(q)q^2\bigr)
\]
rational arithmetic operations, using $\sum_{m\mid q}\varphi(m)=q$.
This is an arithmetic-operation bound, not a constant-cost bit model
or a claim of optimality. The denominator estimate supplies a
conservative size control at positive integer orders.

\section{All-orders Hurwitz--Stieltjes identities}\label{sec:stieltjes}

Write the Laurent expansion as
\begin{equation}\label{eq:stieltjes-laurent}
 \zeta(1+t,x)=\frac1t+
 \sum_{n\ge0}\frac{(-1)^n}{n!}\gamma_n(x)t^n,
 \qquad x>0.
\end{equation}
Here $\gamma_n^{(k)}$ always means $k$ derivatives in $x$; a derivative
on $C(s)$ is in the spectral variable. The coefficient identities in
this section are analytic theorems, distinct from formal rank claims.

\begin{theorem}[Derived Stieltjes reduction]\label{thm:derived-stieltjes}
For $k\ge1$, $n\ge0$, and $x_r=r/q$ with the endpoint convention,
\begin{equation}\label{eq:derived-stieltjes}
 \boxed{\displaystyle
 \gamma_n^{(k)}(x_r)=
 \sum_{u\in U_q}\sum_{j=0}^n\binom nj
 W_{r,u}^{[j]}(k+1)\,
 \gamma_{n-j}^{(k)}(u/q).}
\end{equation}
All the coefficient functions are given by the finite formula
\eqref{eq:C-finite} and its derivatives.
\end{theorem}
\begin{proof}
Differentiating Hurwitz zeta in its parameter yields
\[
 \partial_x^k\zeta(1+t,x)=(-1)^k(1+t)_k\zeta(k+1+t,x).
\]
The pole in \eqref{eq:stieltjes-laurent} is independent of $x$, so for
$k\ge1$ the left side is
$\sum_n(-1)^n\gamma_n^{(k)}(x)t^n/n!$.
Apply \eqref{eq:Hurwitz-C} at spectral parameter $k+1+t$ and multiply
by the common scalar $(-1)^k(1+t)_k$. Taylor multiplication gives
\eqref{eq:derived-stieltjes}, because
$(-1)^j C_{r,u}^{(j)}=W_{r,u}^{[j]}$.
\end{proof}

This proof also explains why the harmonic-number master identity in
Chapter 8 is compatible with the same reduction matrix at every
Stieltjes order: the rising factorial is common to every rational
argument. It does not create new independent distribution coordinates.

\begin{theorem}[Underived Stieltjes reduction, with the pole term]\label{thm:underived}
For $n\ge0$,
\begin{equation}\label{eq:underived-stieltjes}
 \boxed{\displaystyle
 \gamma_n(x_r)=
 \sum_{u\in U_q}\sum_{j=0}^n\binom nj
 W_{r,u}^{[j]}(1)\gamma_{n-j}(u/q)
 -\frac1{n+1}\sum_{u\in U_q}W_{r,u}^{[n+1]}(1).}
\end{equation}
\end{theorem}
\begin{proof}
Apply \eqref{eq:Hurwitz-C} at $s=1+t$. The coefficients are analytic
there. Multiplying their Taylor expansions by
\eqref{eq:stieltjes-laurent}, the coefficient of $t^n$ receives the
usual product terms and also
$\sum_u C_{r,u}^{(n+1)}(1)/(n+1)!$ from the pole $1/t$.
Multiplication by $(-1)^n n!$ converts that additional term to the
negative last term of \eqref{eq:underived-stieltjes}. The coefficient
of $1/t$ agrees because $\sum_u C_{r,u}(1)=1$.
\end{proof}

Omitting the last term would already give an incorrect formula for
$\gamma_0$. The point is not a rank correction: a known scalar right
side does not change a homogeneous rank. It is a correction to the
actual analytic value identity.

\begin{example}[An explicit first-order identity]
From \eqref{eq:q6-one}, direct differentiation at $s=1$ gives
\[
 W_{2,1}^{[1]}(1)=\frac{10}{9}\log2,\qquad
 W_{2,5}^{[1]}(1)=\frac89\log2.
\]
Consequently
\begin{equation}\label{eq:gamma0-six}
 \gamma_0(1/3)=\frac23\gamma_0(1/6)+\frac13\gamma_0(5/6)-2\log2.
\end{equation}
At $s=2$, the analogous coefficients give
\begin{align}\label{eq:gamma1prime-six}
 \gamma_1'(1/3)={}&\frac4{15}\gamma_1'(1/6)
 +\frac1{15}\gamma_1'(5/6)\notag\\
 &+\frac{68\log2}{225}\gamma_0'(1/6)
 +\frac{32\log2}{225}\gamma_0'(5/6).
\end{align}
Since $\gamma_0'=-\psi'$, the signs can also be checked against the
parameter-differentiated digamma multiplication identity.
The verification program checks \eqref{eq:underived-stieltjes} through
$n=2$ and \eqref{eq:gamma1prime-six} by numerical parameter
differentiation of Stieltjes functions.
\end{example}

\section{Cyclotomic polylogarithm traces}\label{sec:traces}

The same finite module applies to a second evaluation. Put
$\zeta_m=\exp(2\pi\ii/m)$ and evaluate
$e_x\mapsto\Li_s(\ee^{2\pi\ii x})$.
For $\Ree s>1$, grouping the defining series by multiples of $p$ gives
\[
 \sum_{py=x}\Li_s(\ee^{2\pi\ii y})
 =p^{1-s}\Li_s(\ee^{2\pi\ii x}).
\]
Thus the weights are now $A_p=p^{1-s}$, and analytic continuation
extends the identity wherever the individual expressions are defined.

For a primitive character $\chi$ of conductor $f\mid q$, define
\begin{equation}\label{eq:trace-definition}
 \cP_{q,\chi}(s)=
 \sum_{u\in U_q}\chi(u)\Li_s(\zeta_q^u),\qquad
 \tau(\chi)=\sum_{a\in U_f}\chi(a)\zeta_f^a.
\end{equation}
At conductor one take $\tau(\mathbf1)=1$ and $L(s,\mathbf1)=\zeta(s)$.

\begin{theorem}[Conductor-lifted trace identity]\label{thm:trace}
As a meromorphic identity, and with removable singularities interpreted
by continuation,
\begin{equation}\label{eq:trace-formula}
 \boxed{\displaystyle
 \cP_{q,\chi}(s)=
 (q/f)^{1-s}
 \prod_{\substack{p\mid q\\p\nmid f}}
 \bigl(1-\chi(p)p^{s-1}\bigr)
 \tau(\chi)L(s,\overline\chi).}
\end{equation}
For every $q>1$, the left side is an entire function of the spectral
variable $s$.
\end{theorem}
\begin{proof}
At primitive conductor $f>1$, absolute convergence and the primitive
Gauss-sum identity give
\[
 \sum_{a\in U_f}\chi(a)\Li_s(\zeta_f^a)
 =\sum_{n\ge1}\frac{\sum_a\chi(a)\zeta_f^{an}}{n^s}
 =\tau(\chi)L(s,\overline\chi).
\]
The Gauss-sum convention, including the complex conjugate on $\chi$,
is the one in~\cite{dlmf-dirichlet}. Applying
\eqref{eq:conductor-normal} with $A_p=p^{1-s}$ multiplies this
primitive trace by $E_{q,\chi}$. Extracting $p^{1-s}$ from each
first-appearance factor gives precisely the factor in
\eqref{eq:trace-formula}. The conductor-one case follows directly
from $\Li_s(1)=\zeta(s)$ and the same polynomial identity.

For completeness, when $z=\zeta_q^u\ne1$, regrouping by residues gives
\begin{equation}\label{eq:polylog-Hurwitz-DFT}
 \Li_s(z)=q^{-s}\sum_{a=1}^q z^a\zeta(s,a/q).
\end{equation}
Hurwitz zeta has only its simple pole at $s=1$, whose coefficient in
this finite sum is $\sum_a z^a=0$. Hence $\Li_s(z)$ is entire in $s$.
All arguments in the primitive level-$q$ trace differ from one.
\end{proof}

\begin{remark}[The character phase is consequential]
For the order-four character modulo five with $\chi(2)=\ii$, the
level-ten lifting factor is $2^{1-s}-\ii$, not
$2^{1-s}+\ii$. A complex-character numerical check is included to
catch an accidental interchange of $\chi$ and $\overline\chi$.
Real quadratic examples alone would not detect that error.
\end{remark}

\subsection{Exact nonprincipal vanishing orders}
For a nonprincipal primitive $\chi$, set
\[
 \rho_\chi=\#\{p\mid q:p\nmid f,\ \chi(p)=1\}.
\]
These are the same conductor statistics that controlled the
primitive-grid jet defects at Hurwitz order zero.

\begin{theorem}[Character trace zeros at order one]\label{thm:trace-zero}
The entire function $\cP_{q,\chi}(s)$ has a zero of exact order
$\rho_\chi$ at $s=1$, where order zero means a nonzero value.
All lower derivatives vanish, and
\begin{align}\label{eq:trace-leading}
 \cP_{q,\chi}^{(\rho_\chi)}(1)
 ={}&(-1)^{\rho_\chi}\rho_\chi!\,
 \tau(\chi)L(1,\overline\chi)\notag\\
 &\times\prod_{\substack{p\mid q,\ p\nmid f\\\chi(p)=1}}\log p
 \prod_{\substack{p\mid q,\ p\nmid f\\\chi(p)\ne1}}(1-\chi(p)).
\end{align}
\end{theorem}
\begin{proof}
Each factor with $\chi(p)=1$ in \eqref{eq:trace-formula} is
$1-p^{s-1}=-(s-1)\log p+O((s-1)^2)$.
All other displayed factors are nonzero at $s=1$.
A primitive Gauss sum is nonzero, and Dirichlet's nonvanishing theorem
states $L(1,\overline\chi)\ne0$ for nonprincipal characters
\cite[Eq.~25.15.9]{dlmf-dirichlet}. Thus there is no additional zero.
Taking the first nonzero Taylor coefficient proves the formula.
\end{proof}

Increasing the exponents of primes already present in $q$ changes
higher Taylor coefficients but not this first nonzero derivative:
$(q/f)^{1-s}$ equals one at $s=1$. This invariance provides an
additional family of identities between traces at different levels.

\subsection{The principal trace and its pole cancellation}
Write $\cP_q=\cP_{q,\mathbf1}$ and define
\[
 r=\omega(q),\quad P=\prod_{p\mid q}\log p,\qquad
 B=\log q-\tfrac12\log\rad(q).
\]
The conductor-one factor contains a pole, so the vanishing order is
one less than the number of primes, not equal to that number.

\begin{theorem}[Principal trace jet generator]\label{thm:principal}
In a neighborhood of $t=0$,
\begin{equation}\label{eq:principal-generator}
 \boxed{\displaystyle
 \cP_q(1+t)=(-1)^rP\,t^{r-1}\,
 \ee^{-Bt}\,[t\zeta(1+t)]
 \prod_{p\mid q}\frac{\sinh(t\log p/2)}{t\log p/2}.}
\end{equation}
In particular its zero at $s=1$ has exact order $r-1$, and
\begin{align}
 \cP_q^{(j)}(1)&=0 &&(0\le j<r-1),\label{eq:principal-vanish}\\
 \cP_q^{(r-1)}(1)&=(-1)^r(r-1)!P,\label{eq:principal-first}\\
 \cP_q^{(r)}(1)&=(-1)^r r!P(\gamma-B),\label{eq:principal-second}\\
 \cP_q^{(r+1)}(1)&=(-1)^r(r+1)!P
 \left[-\gamma_1-B\gamma+\frac{B^2}{2}
 +\frac1{24}\sum_{p\mid q}(\log p)^2\right].
 \label{eq:principal-third}
\end{align}
Here $\gamma_1$ is the ordinary first Stieltjes constant in
\eqref{eq:stieltjes-laurent}.
\end{theorem}
\begin{proof}
The principal case of \eqref{eq:trace-formula} is
\[
 \cP_q(1+t)=q^{-t}\prod_{p\mid q}(1-p^t)\zeta(1+t).
\]
Use $1-\ee^x=-x\ee^{x/2}\sinh(x/2)/(x/2)$ in every factor.
This proves \eqref{eq:principal-generator}.
The factors in square brackets and in the hyperbolic-sine product are
analytic and equal to one at zero. Expanding them gives
\[
 t\zeta(1+t)=1+\gamma t-\gamma_1t^2+O(t^3),\qquad
 \prod_p\frac{\sinh(t\log p/2)}{t\log p/2}
 =1+\frac{t^2}{24}\sum_p(\log p)^2+O(t^4).
\]
Multiplication by $\ee^{-Bt}$ gives all the displayed derivatives.
\end{proof}

\subsection{Concrete exact identities}
The following examples are proved consequences of the preceding
theorems, not integer-relation conjectures. The derivatives are in the
\emph{order} of the polylogarithm, not its argument.

At level thirty,
\begin{align}
 \sum_{u\in U_{30}}\Li_1(\zeta_{30}^u)&=0,\label{eq:q30-zero}\\
 \sum_{u\in U_{30}}\Liord{1}{\zeta_{30}^u}&=0,\label{eq:q30-first}\\
 \sum_{u\in U_{30}}\Liord{2}{\zeta_{30}^u}
 &=\boxed{-2\log2\log3\log5}.\label{eq:q30-second}
\end{align}
The third derivative of this trace is
$-6\log2\log3\log5\,[\gamma-\frac12\log30]$.
At level $210$, derivatives of orders zero, one, and two vanish, while
\begin{equation}\label{eq:q210}
 \sum_{u\in U_{210}}\Liord{3}{\zeta_{210}^u}
 =6\log2\log3\log5\log7.
\end{equation}

For the quadratic character $\chi_3$ with $\chi_3(1)=1$ and
$\chi_3(2)=-1$, the primitive trace at order one is
$\tau(\chi_3)L(1,\chi_3)=\ii\pi/3$. Since
$\chi_3(7)=\chi_3(13)=1$, one obtains
\begin{align}
 \sum_{u\in U_{21}}\chi_3(u)\Liord{1}{\zeta_{21}^u}
 &=-\frac{\ii\pi}{3}\log7,\label{eq:chi3-21}\\
 \sum_{u\in U_{273}}\chi_3(u)\Liord{2}{\zeta_{273}^u}
 &=\frac{2\ii\pi}{3}\log7\log13.\label{eq:chi3-273}
\end{align}
In the first equation the order-zero trace vanishes; in the second,
both order-zero and first-order traces vanish.
Likewise the character $\chi_4$ gives
\begin{equation}\label{eq:chi4-20}
 \sum_{u\in U_{20}}\chi_4(u)\Liord{1}{\zeta_{20}^u}
 =-\frac{\ii\pi}{2}\log5.
\end{equation}
These primitive trace constants follow directly from
$\Li_1(z)=-\log(1-z)$ with the values continued from inside the unit
disk, so the signs do not depend on an unspecified logarithm branch.

For the real quadratic character $\chi_5$, positive on residues
$1,4$ and negative on $2,3$, let $\phi=(1+\sqrt5)/2$.
Pairing conjugates in the primitive trace gives
\[
 \cP_{5,\chi_5}(1)
 =2\log\frac{\sin(2\pi/5)}{\sin(\pi/5)}=2\log\phi.
\]
As $\chi_5(11)=1$,
\begin{equation}\label{eq:chi5-55}
 \sum_{u\in U_{55}}\chi_5(u)\Liord{1}{\zeta_{55}^u}
 =-2\log\phi\log11.
\end{equation}
This derivation needs no numerical integer-relation search or
unproved independence statement about logarithms.

\subsection{A stable independent formula for numerical checks}
Differentiation of \eqref{eq:polylog-Hurwitz-DFT} at its cancelled pole
gives, for $z=\zeta_q^u\ne1$,
\begin{equation}\label{eq:stable-polylog-jet}
 \Liord{k}{z}
 =\frac{(-1)^k}{q}\sum_{a=1}^q z^a
 \sum_{j=0}^k\binom kj(\log q)^{k-j}\gamma_j(a/q).
\end{equation}
There is no pole term, because its complete Fourier coefficient
vanishes identically. This supplies a finite Stieltjes-coordinate
evaluator for the trace identities, independently of their
Euler-factor right sides. In development, direct tiny-step numerical
differentiation of a polylogarithm routine near integral order gave
an unreliable trace derivative; the pole-cancelled formula avoids
that cancellation problem. This observation is an implementation
diagnostic, not a claimed theorem about all versions of a library.

\section{The Gaussian \texorpdfstring{$S_4$}{S4} conjecture: an exact obstruction}\label{sec:s4}

The manuscript uses decreasing nested indices:
\[
 \Li_{a,b}(x,y)=\sum_{m>n\ge1}\frac{x^my^n}{m^a n^b},\qquad
 g_{a,b}=\Imm\Li_{a,b}(\ii,1),\quad G=\beta(2).
\]
Its equation \code{gauss:eq:S4-closed} is the following candidate.

\begin{conjecture}[The manuscript's weight-five evaluation]\label{conj:S4}
With $H_0=0$ and $H_n=\sum_{j=1}^n1/j$,
\begin{equation}\label{eq:S4-conjecture}
 S_4:=\sum_{n\ge0}\frac{(-1)^nH_n}{(2n+1)^4}
 =\frac{4g_{4,1}-3g_{3,2}-9g_{2,3}}7
 +\frac{\pi^5}{224}-\frac{27G\zeta(3)}{224}
 -2\beta(4)\log2.
\end{equation}
\end{conjecture}

We retain this as a conjecture here. The new rigorous result is a
specific obstruction to deriving it from products of two single
polylogarithms alone, in the exact presentation defined next. This is
more informative than a failed numerical search, but it must not be
misread as an obstruction to all possible proofs.

\begin{lemma}[An exact even-inner-index filter]\label{lem:S4-filter}
One has
\begin{equation}\label{eq:S4-filter}
 S_4=\Imm\bigl(\Li_{4,1}(\ii,1)+\Li_{4,1}(\ii,-1)\bigr).
\end{equation}
\end{lemma}
\begin{proof}
The inner sum in the right side is
\[
 \sum_{n<m}\frac{1+(-1)^n}{n}=H_{\lfloor(m-1)/2\rfloor}.
\]
The imaginary part of $\ii^m$ vanishes at even $m$ and is $(-1)^j$
at $m=2j+1$. Substitution gives the defining sum for $S_4$.
The double sums are absolutely convergent, so this rearrangement is
legitimate.
\end{proof}

\subsection{A fully specified finite law vocabulary}
Introduce formal coordinates
\[
 X_{a;r,t}=\Imm\Li_{a,5-a}(\ii^r,\ii^t),\qquad 1\le a\le4.
\]
Complex conjugation identifies $X_{a;-r,-t}=-X_{a;r,t}$, and a pair
of even residues has imaginary part zero. Choose the six representatives
\begin{equation}\label{eq:S4-pairs}
 (r,t)\in\{(0,1),(1,0),(1,1),(1,2),(1,3),(2,1)\}.
\end{equation}
They give $24$ double coordinates. Append three formal single-value
product coordinates
\[
 A=\pi^5,\qquad B=G\zeta(3),\qquad C=\beta(4)\log2.
\]
For each $p+v=5$, $p,v\ge1$, and $r,t\in\Z/4\Z$, include the
imaginary parts of the two identities
\begin{align}
 \Li_p(\ii^r)\Li_v(\ii^t)
 ={}&\Li_{p,v}(\ii^r,\ii^t)+\Li_{v,p}(\ii^t,\ii^r)
 +\Li_5(\ii^{r+t}),\label{eq:S4-stuffle}\\
 \Li_p(\ii^r)\Li_v(\ii^t)
 ={}&\sum_{j=0}^{p-1}\binom{v-1+j}{j}
 \Li_{v+j,p-j}(\ii^t,\ii^{r-t})\notag\\
 &+\sum_{j=0}^{v-1}\binom{p-1+j}{j}
 \Li_{p+j,v-j}(\ii^r,\ii^{t-r}).\label{eq:S4-shuffle}
\end{align}
The first is the colored stuffle law. The second follows by shuffling
$0^{p-1}(\ii^r)^{-1}$ with $0^{v-1}(\ii^t)^{-1}$ in the
outer-letter-first iterated-integral convention; the binomial
coefficients count the interspersed zero letters.

Use the standard single-value reductions at fourth roots of unity.
Their imaginary products at total weight five lie in the three atoms
$A,B,C$. At the divergent factor $\Li_1(1)$, use regularized constant
term zero. Only a linear divergence polynomial occurs here: a weight-one
factor can multiply a convergent weight-four factor, but two divergent
weight-one factors cannot occur at total weight five in this vocabulary.
Thus shuffle and stuffle regularization agree on the degree-zero and
degree-one terms being used; the general regularization framework is
standard~\cite{racinet}. The formal row-space result below would remain
well-defined even if these boundary rows were omitted; we report that
smaller matrix separately.

After conjugation reduction and removal of zero rows there are $96$
rows and $27$ columns. Duplicate nonzero rows are intentionally retained,
so the law-generation specification is simple and replayable. Without
products containing $\Li_1(1)$ there are $88$ rows.

\begin{theorem}[A finite row-space obstruction]\label{thm:S4-obstruction}
Let $M$ be the $96\times27$ rational matrix just defined. Its rank is
$20$. The row corresponding, via \eqref{eq:S4-filter}, to
$224$ times the difference between the two sides of
\eqref{eq:S4-conjecture} is
\begin{equation}\label{eq:S4-target}
 T=96X_{4;1,0}+96X_{3;1,0}+288X_{2;1,0}
 +224X_{4;1,2}-A+27B+448C.
\end{equation}
It does not lie in the row span of $M$. The augmented matrix has rank
$21$. The unregularized $88$-row matrix has rank $19$.
\end{theorem}
\begin{proof}
The nonmembership assertion has a particularly small certificate.
Define a linear functional $w$ on the $27$ coordinates by the following
eight nonzero entries, and set it to zero on every other coordinate,
including $A,B,C$:
\begin{center}
\begin{tabular}{cr@{\qquad}cr}
\toprule
Coordinate & $w$ & Coordinate & $w$\\\midrule
$X_{1;1,0}$&$-2$&$X_{2;0,1}$&$1$\\
$X_{2;1,0}$&$3$&$X_{2;1,3}$&$-1$\\
$X_{3;0,1}$&$-3$&$X_{3;1,0}$&$-1$\\
$X_{3;1,3}$&$-1$&$X_{4;0,1}$&$2$\\
\bottomrule
\end{tabular}
\end{center}
Substitution in the two finite row formulas
\eqref{eq:S4-stuffle}--\eqref{eq:S4-shuffle}, for their stated index
ranges, gives $Mw=0$. This is $96$ exact rational zero checks, recorded
in the certificate. On the target, however,
\[
 w(T)=96(-1)+288(3)=768\ne0.
\]
Thus no linear combination of the declared rows equals $T$.
The rank assertions are additional finite exact calculations: Gaussian
elimination gives ranks $20,21,19$, respectively. Both the generator
using symbolic rational arithmetic and an independent standard-library
rational replay program are supplied. The nonmembership proof itself
needs only the displayed eight-entry witness, not the rank calculation.
\end{proof}

This theorem remains true after adding arbitrary relations supported
only on the three product atoms, because $w$ annihilates those atoms.
It does not remain a theorem about an enlarged multiple-polylogarithm
presentation unless the new rows are also checked. Products involving
doubles, elimination through depth three, distribution at other levels,
or functional equations at transformed arguments may provide the missing
relation. The witness directs future searches toward precisely the
relations on which it does not vanish.

\subsection{Fresh numerical evidence, with its exact limitation}
An independent integral comparison uses
\begin{align}
 S_4&=-\frac16\int_0^1
 \frac{(-\log x)^3\log(1+x^2)}{1+x^2}\,dx,\label{eq:S4-integral}\\
 \Li_{a,b}(z,1)&=\frac1{(a-1)!}\int_0^1
 (-\log t)^{a-1}\frac{z\Li_b(zt)}{1-zt}\,dt.
 \label{eq:double-integral}
\end{align}
For \eqref{eq:S4-integral}, expand
$\sum_{n\ge0}H_n(-x^2)^n=-\log(1+x^2)/(1+x^2)$ and integrate.
For \eqref{eq:double-integral}, interchange absolutely convergent
power series inside the integral and use
$\int_0^1t^{m-1}(-\log t)^{a-1}\,dt=(a-1)!/m^a$.
The relevant boundary values follow by dominated convergence.

At $55$ decimal working digits the independent integral evaluations give
\[
 S_4=-0.01049345629733504345335675100620175196351061887731976151\ldots
\]
and the residual in \eqref{eq:S4-conjecture} is approximately
$7.17\times10^{-57}$. This is floating-point agreement, not an interval
proof and not a proof that the residual is exactly zero. The claim that
has been proved is the finite-presentation obstruction, not the
conjectured evaluation.

\section{Corrections and clarifications for the manuscript}\label{sec:audit}

\subsection{Replace the distribution-rank experiment by a defined theorem}
Replace the experimental-only formulation at
\code{tower:thm:rank} by Corollary~\ref{cor:anchored}, with an explicit
prime/divisor row definition. Retain the explanation that formal quotient
dimension does not prove numerical independence. The same addition
supplies the uniform parity count in Chapter 7 and all-order jet
extensions in Chapter 8. The historical data may remain as provenance,
but should not be described as the proof of the uniform statement.

At exceptional weights, do not infer a rank jump from a failed inverse.
Theorem~\ref{thm:universal} gives a globally valid conductor basis;
Theorems~\ref{thm:determinant} and \ref{thm:smith} identify the separate
failure of primitive-grid coordinates. A proof based only on their
inversion would miss this distinction.

\subsection{The cubic log-gamma moment already has a named-constant formula}
Chapter 7, \code{integral:neg:cubic}, correctly observes that a small
basket of single zeta derivatives did not explain the cubic moment.
However, its claim that finding the appropriate named Tornheim-type
atoms is still open is superseded by Bailey--Borwein--Borwein:
their paper explicitly evaluates the cubic moment in Theorem 6 and
equation (11), using derivatives of Tornheim--Witten zeta
functions~\cite{bbb}. The appropriate correction is to distinguish an
existing formula in that larger vocabulary from the separate question
of reduction to a smaller vocabulary. Appendix~\ref{app:cubic} gives
the formula in the notation used here and an independent Fourier
coefficient derivation.

This is a literature correction, not a new evaluation claimed by this
article. It does not imply that all derivatives appearing there reduce
to ordinary zeta derivatives or classical single polylogarithms.

\subsection{At mixed Gaussian points the relevant involution swaps arguments}
The stuffle relation identifies a symmetric combination under
\[
 (a,b;x,y)\longmapsto(b,a;y,x),
\]
not merely $(a,b)\mapsto(b,a)$ with $(x,y)$ fixed. Thus at a mixed
point the complementary antisymmetric combination is
\[
 \Li_{a,b}(x,y)-\Li_{b,a}(y,x).
\]
It need not vanish when $a=b$ and $x\ne y$. The manuscript's
``all depth-two content'' sentence following its symmetric Gaussian
formulas should be restricted to the paired involution. Nothing in
this correction asserts that a nonzero antisymmetric combination is
numerically independent of single-value products.

\begin{proposition}[A nonzero mixed diagonal]\label{prop:mixed-diagonal}
One has the strict inequality
\begin{equation}\label{eq:mixed-positive}
 \Imm\bigl(\Li_{2,2}(-1,\ii)-\Li_{2,2}(\ii,-1)\bigr)>0.
\end{equation}
\end{proposition}
\begin{proof}
The general double-integral formula, with $\Li_b(xyt)$ in the
numerator, gives their difference as
\[
 -\int_0^1(-\log t)\,
 \frac{(1+\ii)\Li_2(-\ii t)}{(1+t)(1-\ii t)}\,dt.
\]
Write $\operatorname{Ti}_2(t)=\int_0^t\arctan(v)\,dv/v>0$ for
$0<t\le1$. Then
$\Li_2(-\ii t)=\tfrac14\Li_2(-t^2)-\ii\operatorname{Ti}_2(t)$.
The imaginary part of the displayed integrand, including $-\log t$,
is
\[
 \frac{-\log t}{(1+t)(1+t^2)}
 \left[(1-t)\operatorname{Ti}_2(t)
       -\frac{1+t}{4}\Li_2(-t^2)\right].
\]
It is strictly positive on $(0,1)$, since $\Li_2(-t^2)<0$.
The integral converges, proving \eqref{eq:mixed-positive}.
\end{proof}

\subsection{Keep three meanings of a certificate distinct}
The polynomial row certificates prove formal identities from a defined
law vocabulary. Their analytic evaluations require the convergence and
continuation arguments supplied above. Floating-point residuals test
implementations but prove neither formal row membership nor equality of
general transcendental numbers. In particular, Theorem~\ref{thm:S4-obstruction}
and the small numerical residual for $S_4$ are fully compatible facts.
A proof-assistant port would add another layer of verification; none is
claimed for this package.

\section{Reproducible computations and integration}\label{sec:computations}

\subsection{Exact checks actually executed}
The package's exact verification program tests every $2\le q\le60$.
For each denominator, it uses five weight patterns: $A_p=p^2$, all
$A_p=1$, all $A_p=0$, all $A_p=-1$, and a mixed zero/nonzero pattern.
In each of the resulting $295$ cases it checks the full rank,
anchored rank, and applicable even/odd ranks with exact rational
arithmetic. These include deliberately resonant and zero weights.

A second set of $295$ cases uses spectral orders
$s=-2,-1,1,2,3$ for the same denominators. It constructs
\eqref{eq:C-finite} exactly and checks every divisor row, the identity
primitive submatrix, the predicted support counts, the sign rule, and
the endpoint formula. Negative orders use finite rational continuation,
not an infinite series. Four complete matrices are included as JSON
certificates, for $(q,s)=(12,2),(21,-1),(30,3),(60,2)$.

The character verification covers the same denominator range. Characters
are represented by rational phases modulo one, not approximate complex
numbers. It checks conductor counts, the root multiplicities and pure
powers in \eqref{eq:determinant}, and the central primitive-jet rank
profiles. Character examples for levels $12,21,30$ are supplied.

The $S_4$ matrix and the eight-entry obstruction witness are regenerated
from the declared laws, then replayed independently with
\code{fractions.Fraction}. The replay checks all $96$ annihilations,
the target pairing $768$, and ranks $20$ and $21$.
Finite regression ranges support the software; the uniform statements
are established by the proofs, not by extrapolation from these tests.

\subsection{Numerical diagnostics actually executed}
The numerical program was run at $55$ decimal working digits. It checks
seven trace/phase identities, three underived Stieltjes identities, one
parameter-differentiated Stieltjes identity, and the $S_4$ integral
comparison. Trace derivatives use \eqref{eq:stable-polylog-jet} with
additional guard digits, while a nonintegral-order complex-character
phase check uses direct polylogarithms. No target Euler-factor identity
is used to evaluate its own trace side.

\begin{center}
\begin{tabular}{p{0.56\textwidth}r}
\toprule
Diagnostic family & Largest observed absolute residual\\\midrule
Trace and complex-character checks & $5.77\times10^{-56}$\\
Underived Stieltjes checks, $n=0,1,2$ & $8.16\times10^{-56}$\\
Derived Stieltjes check, $k=n=1$ & $3.27\times10^{-55}$\\
$S_4$ independent integral comparison & $7.17\times10^{-57}$\\
\bottomrule
\end{tabular}
\end{center}

Residuals smaller than the working precision can occur because rounded
real parts coincide or because extra guard digits are retained. They
are not claims of that many correct decimal places. These experiments
use arbitrary-precision floating point, not interval arithmetic, and
the JSON records say so explicitly.

\subsection{Files and reproduction}
The article is self-contained in \code{article/distribution\_jets.tex}.
The principal programs are:
\begin{center}
\begin{tabular}{p{0.43\textwidth}p{0.50\textwidth}}
\toprule
File under \code{code/} & Purpose\\\midrule
\code{distribution.py} & Exact distribution rows, rational ranks, finite normal forms.\\
\code{verify\_exact.py} & Uniform-range regression checks and matrix certificates.\\
\code{verify\_characters.py} & Exact character phases, determinant factors, and jet profiles.\\
\code{verify\_s4\_presentation.py} & Generate the declared Gaussian row system and witness.\\
\code{replay\_certificate.py} & Standard-library replay of the $S_4$ certificate.\\
\code{verify\_numeric.py} & Separate trace, Stieltjes, and $S_4$ numerical diagnostics.\\
\bottomrule
\end{tabular}
\end{center}

The distribution and replay modules use only the Python standard library.
The character generator and symbolic Gaussian generator use SymPy;
numerical diagnostics use mpmath. The versions used in this run are
recorded in the package. The \code{integration/} directory provides
editorial replacement text and a label map rather than silently
rewriting the live manuscript. Source paths, inspected Git blob hashes,
and the repository commit are recorded in \code{provenance.json}.
No repository file was changed remotely.

\section{Further research questions}\label{sec:future}

\subsection{Complete the proof of the weight-five conjecture}
The most direct next problem is to prove or refute
Conjecture~\ref{conj:S4}. The exact witness makes an enlarged search
more targeted: generate candidate new laws, apply $w$ first, and
prioritize laws with nonzero pairing. Promising sources include
products of a single and a double polylogarithm with elimination of
triples, regularized distribution relations, and argument-changing
functional equations. A successful proof should include the additional
law instances and an exact final row combination, not just a larger
integer-relation search.

\subsection{Integral and modular versions of the weighted module}
The proof of Theorem~\ref{thm:universal} divides by character group
orders and works over a characteristic-zero splitting field. What is
the exact integral structure of the independently weighted module?
At ordinary weights, classical integral freeness provides an important
benchmark. At general weights, one should determine torsion and the
primes at which the conductor Fourier method ceases to describe the
integral lattice. In positive characteristic, nonsemisimple character
actions can change the argument; the characteristic-zero theorem
alone does not settle that problem.

\subsection{Multivariate jets and intersection multiplicities}
Independent variables $t_p$ in $A_p=\exp(t_p)$ give a multivariate
version of the normal form. The module remains free by base change,
but the primitive-grid inclusion has a union of coordinate and
cyclotomic divisors rather than one-variable Smith factors. It would
be useful to compute its Fitting ideals, truncated multigraded ranks,
and optimal minimal jet data. The determinant gives the first
invariant, not the entire multivariate classification.

\subsection{Sparse algorithms and certified numerical evaluation}
The residue-subgroup support theorem can be used to avoid dense
matrices. Can one give output-sensitive complexity bounds in terms of
$\sum_r |\operatorname{supp}C_r|$, rather than $q^2$? Can the
conductor basis be combined with interval arithmetic to evaluate
spectral jets uniformly near resonances, without catastrophic
cancellation? The coefficients have explicit denominators at positive
integer orders; a certified remainder analysis for analytic jets would
complement that exact arithmetic structure.

\subsection{Trace identities at greater depth}
For multiple polylogarithms, distributing a root argument interacts
with several nested indices and with prefix products. A useful
extension would identify an explicit finite module whose degree-one
case is the present theory and whose higher-depth relations remain
compatible with shuffle and stuffle. One should not assume that the
rank simply becomes a power of $\varphi(q)$: depth, convergence,
regularization, and diagonal terms all enter.

\subsection{Arithmetic meaning of the trace-jet constants}
Theorems~\ref{thm:trace-zero} and \ref{thm:principal} identify exact
vanishing orders and first nonzero coefficients. The higher
coefficients involve $L$-derivatives, Stieltjes constants, and
logarithms. Which additional functional equations reduce their
arithmetic vocabulary? Any claim of minimality must specify a
comparison algebra and distinguish formal, motivic, and numerical
independence. The present rank theorem alone gives none of those
numerical lower bounds.

\subsection{Log-gamma moments beyond the known Tornheim formula}
The cubic moment should be studied starting from its existing
Tornheim--Witten derivative expression, not by rediscovering the named
atoms. Useful concrete questions include reductions among the
particular derivatives $T_c,T_{ac},T_{ab}$, better certified evaluation
of those derivatives, and the organization of higher log-gamma moments
by Fourier resonance patterns. Any proposed smaller closed form should
be tested against the established larger-vocabulary identity.

\subsection{A proof-assistant development}
A staged formalization can begin with the finite module: define the
fiber rows, prove prime sufficiency, establish character-block edges,
and verify the monic conductor normal form. This isolates finite
algebra from analytic continuation. The second stage can formalize
Hurwitz and polylogarithm evaluations and their local coefficient
identities. The small $S_4$ annihilator is a separate attractive target
for a finite rational certificate checker. None of these future
formalizations is asserted to have been completed by the present work.

\clearpage
\appendix
\section{The known cubic log-gamma formula in Tornheim notation}\label{app:cubic}

Define, in an open neighborhood of $(1,1,1)$,
\[
 T(a,b,c)=\sum_{m,n\ge1}\frac1{m^a n^b(m+n)^c}.
\]
The series and all finite partial derivatives there converge
absolutely. Let $T_c,T_{ac},T_{ab}$ mean the indicated partial
derivatives evaluated at $(1,1,1)$, and put
$L=\log(2\pi)$, $A=\gamma+L$. A form of the
Bailey--Borwein--Borwein evaluation~\cite[Eq.~(11), Theorem~6]{bbb} is
\begin{align}\label{eq:cubic-known}
 \int_0^1\log^3\Gamma(x)\,dx
 ={}&\frac{L^3}{8}+\frac{L\pi^2}{32}+\frac{LA^2}{8}
 -\frac{3LA\zeta'(2)}{2\pi^2}
 +\frac{3L\zeta''(2)}{4\pi^2}+\frac{3\zeta(3)}{16}\notag\\
 &+\frac3{8\pi^2}
 \bigl[2A^2\zeta(3)-2AT_c+2T_{ac}-T_{ab}\bigr].
\end{align}
Here $T_c=2\partial_u\zeta(u,1)|_{u=2}$, where in this sentence
$\zeta(u,1)=\sum_{m>n\ge1}m^{-u}n^{-1}$ is the \emph{double zeta},
not the Hurwitz zeta. This follows by partial fractions and identifies
our notation with the published formula.

A direct derivation is short. Kummer's Fourier expansion gives the
mean $L/2$ and positive-frequency coefficients
\[
 \alpha_n=\frac1{4n}-\frac{\ii(A+\log n)}{2\pi n}
 \qquad(n\ge1)
\]
for $\log\Gamma(x)$. The negative-frequency coefficients are their
complex conjugates. Taking the zero-frequency coefficient of the
cube yields
\begin{equation}\label{eq:cubic-fourier}
 (L/2)^3+6(L/2)\sum_{n\ge1}|\alpha_n|^2
 +6\Ree\sum_{m,n\ge1}\alpha_m\alpha_n
 \overline{\alpha_{m+n}}.
\end{equation}
The first sum is
\[
 \sum_{n\ge1}|\alpha_n|^2
 =\frac{\pi^2}{96}+\frac{A^2}{24}
 -\frac{A\zeta'(2)}{2\pi^2}+\frac{\zeta''(2)}{4\pi^2}.
\]
For the double sum, expand the three factors. The constant term uses
$T(1,1,1)=2\zeta(3)$. Terms linear in the logarithms leave only the
$\log(m+n)$ sum, namely $-T_c$; the quadratic terms give
$2T_{ac}-T_{ab}$. Substitution in \eqref{eq:cubic-fourier} gives
\eqref{eq:cubic-known}.

To justify the multiplication, first multiply every Fourier
coefficient by $\rho^{|n|}$ with $0<\rho<1$. The smoothed Fourier
series is absolutely convergent. Since $\log\Gamma$ belongs to
$L^3(0,1)$, these Poisson-smoothed functions converge in $L^3$, so
the integrals of their cubes converge. The logarithmically weighted
Tornheim sums are absolutely convergent near $(1,1,1)$ and dominate
the corresponding smoothed double sums. Passage to the limit in both
representations is therefore legitimate. This derivation is included
to make the proposed correction mathematically self-contained, not to
claim priority for the formula.

\clearpage
\section{Provenance, claim boundaries, and artifact checks}\label{app:provenance}

The research baseline is repository commit
\begin{center}
\code{29d9344771b41df39ff9b7890c76b1f152572dde}.
\end{center}
The source root is
\begin{center}
\code{Analysis/Polylogarithms/docs/manuscript/}.
\end{center}
Relevant inspected blob hashes are recorded in the package's machine-readable
provenance file. In particular, the weight-five conjecture was checked
against the pinned Chapter 4 blob
\code{6171912ff7724784cabb789a066c61223ab298bc}.
The article is a companion contribution: it does not claim to audit
all eleven chapters or every historical identity store.

\begin{center}
\begin{tabular}{p{0.51\textwidth}p{0.41\textwidth}}
\toprule
Claim & Status in this contribution\\\midrule
Universal weighted rank and conductor normal form & Proved over a characteristic-zero splitting field.\\
Primitive-grid determinant and local jet ranks & Proved; exact character checks included.\\
Hurwitz--Stieltjes and cyclotomic trace formulas & Analytic identities proved with stated conventions.\\
$S_4$ evaluation & Retained as conjectural; fresh numerical agreement only.\\
$S_4$ nonmembership in the declared row span & Proved by a finite exact annihilator.\\
Mixed Gaussian diagonal antisymmetry & A strictly positive imaginary difference is proved.\\
Cubic log-gamma evaluation & Existing published result; notation and status corrected.\\
Numerical independence; minimal transcendental basis & Not claimed.\\
Proof-assistant verification; external peer review & Not claimed.\\
\bottomrule
\end{tabular}
\end{center}

The bundled tests separate exact arithmetic from floating-point
experiments. Every certificate has a declared meaning and an explicit
coordinate convention. A checksum manifest is provided for the package
contents. The build and test scripts operate locally and do not write
to the remote repository.

\clearpage
\begin{thebibliography}{99}
\bibitem{repo}
V. Reshetnikov, \emph{Polylogarithms and their Arithmetic Bridges},
ProveIt repository, consolidated manuscript, October 2026.
Inspected commit \code{29d9344771b41df39ff9b7890c76b1f152572dde}.
\url{https://github.com/VladimirReshetnikov/ProveIt/tree/29d9344771b41df39ff9b7890c76b1f152572dde/Analysis/Polylogarithms/docs/manuscript}.

\bibitem{ouyang}
Y. Ouyang, \emph{Group cohomology of the universal ordinary distribution},
arXiv:math/9911032 (1999), especially Proposition 3.1.
\url{https://arxiv.org/abs/math/9911032}.

\bibitem{kubert}
D. S. Kubert, \emph{The $\Z/2\Z$ cohomology of the universal ordinary
distribution}, Bulletin de la Soci\'et\'e Math\'ematique de France
\textbf{107} (1979), 203--224.
\url{https://www.numdam.org/article/BSMF_1979__107__203_0.pdf}.

\bibitem{dlmf-hurwitz}
NIST Digital Library of Mathematical Functions, \S25.11,
\emph{Hurwitz Zeta Function}. Accessed 9 October 2026.
\url{https://dlmf.nist.gov/25.11}.

\bibitem{dlmf-polylog}
NIST Digital Library of Mathematical Functions, \S25.12,
\emph{Polylogarithms}. Accessed 9 October 2026.
\url{https://dlmf.nist.gov/25.12}.

\bibitem{dlmf-dirichlet}
NIST Digital Library of Mathematical Functions, \S25.15,
\emph{Dirichlet $L$-functions}, including (25.15.4), (25.15.6),
and the nonvanishing statement (25.15.9). Accessed 9 October 2026.
\url{https://dlmf.nist.gov/25.15}.

\bibitem{racinet}
G. Racinet, \emph{Doubles m\'elanges des polylogarithmes multiples aux
racines de l'unit\'e}, arXiv:math/0202142 (2002).
\url{https://arxiv.org/abs/math/0202142}.

\bibitem{bbb}
D. H. Bailey, D. Borwein, and J. M. Borwein,
\emph{On Eulerian Log-Gamma Integrals and Tornheim--Witten Zeta Functions},
preprint dated 27 June 2012; Ramanujan Journal \textbf{36} (2015), 43--68.
DOI: 10.1007/s11139-012-9427-1.
\url{https://carmamaths.org/resources/jon/log-gamma.pdf}.
\end{thebibliography}
\end{document}
